\scanpage{65}
% block: warning continued and exception
et même dans le cas particulier ou\inserted{\handwritten{de plus\begin{tikzpicture}[x=1cm,y=1cm,baseline=-.1cm]\draw[thin] (0,.3).. controls (0,.03) and (-.65,.02)..(-.8,-.65);\end{tikzpicture}}} $\handwritten{\mathcal{M}}$ est la catégorie des représentations d'un groupe fini (? ? à vérifier dans littérature qu'un groupe fini n'est pas déterminé par la connaissance de la table de multiplication de \deleted{ses \unsure{lt|it|il}} son ensemble de caractères, et les rangs de ceux-ci ...) Comme exception à cette remarque, signalons cependant le cas où $\handwritten{\mathcal{M}}$ est tannakienne à lien diagonalisable (a\unsure{\deleted{n}l|\deleted{l}n}ors $\handwritten{\Sigma}$ devient un sous-groupe de $K^{\unsure{\times|x|\mathbb{x}|*}}$, \deleted{donc}\inserted{qui} détermine le lien, et la donnée de 3) donne l'homomorphisme $\handwritten{\Sigma\longrightarrow}\operatorname{Br}(k)$ qui suffit à tout déterminer). Regarder le cas où le lien est un groupe de t.m.

% block: c multiplicative type
c) Cas d'un lien qui est de t.m. On suppose connu le caractère numérique \deleted{(fon\unsure{ç|c}} de $\handwritten{\mathcal{M}}$. Prenant le changement de base par une clôture séparable $k'$ de $k$, on trouve que $\handwritten{\Sigma'=\Sigma(\mathcal{M}}_{k'}\deleted{\unsure{x|\times}}\handwritten{)}$ (canoniquement déduit de la donnée \unsure{\deleted{d}n|\deleted{n}d}umérique, avec l'opération de $\handwritten{\pi}=\operatorname{Gal}(k'/k)$ dessus) est muni d'une structure de groupe induite par la multiplication de $K(\handwritten{\mathcal{M}}_{k'})=\underline{Z}^{(\handwritten{\Sigma'})}$, \deleted{et a\unsure{v|u}} respectée par l'action de $\handwritten{\pi}$. Ceci nous donne le lien $G=D(\handwritten{\Gamma})$ (où $\handwritten{\Gamma}$ est le groupe $\handwritten{\Sigma'}$ avec l'action de $\handwritten{\pi}$). L'ensemble $\handwritten{\Sigma}$ s'identifie à l'ensemble des orbites $\handwritten{\sigma}$ de $\handwritten{\pi}$ opérant sur $\handwritten{\Gamma}$. Si $\handwritten{\sigma}$ est une telle orbite, le corps $Z_{\handwritten{\sigma}}$ n'est autre que le corps des invariants $Z_{\handwritten{\sigma}}\handwritten{\subset}k'$ du stabilisateur d'un point de l'orbite $\handwritten{\sigma}$. Si $M_{\handwritten{\sigma}}$ est un objet de type $\handwritten{\sigma}$, on a un homomorphisme
\[
G\handwritten{\longrightarrow}\underline{\operatorname{Aut}}(M_{\handwritten{\sigma}})=\underline{K}_{\handwritten{\sigma}}^\times
\]
(les soulignés indiquent qu'on prend des schémas sur $k$ ...) qui se factorise par le centre $Z_{\handwritten{\sigma}}$ de $K_{\handwritten{\sigma}}$ en
\[
u_{\handwritten{\sigma}}:\ G\handwritten{\unsure{\mathrel{\begin{tikzpicture}[x=1cm,y=1cm,baseline=-.05cm]\draw[thin] (0,0)--(.55,.02)--(.52,.12);\end{tikzpicture}}|\longrightarrow}}\underline{Z}_{\handwritten{\sigma}}^\times\handwritten{\unsure{=|\simeq}}{}_{Z}\handwritten{\prod}_{/k}\underline{G}_m\ ,
\]
et on vérifie aussitôt qu'on a
\[
u_{\handwritten{\sigma}}(\handwritten{\xi})=\handwritten{\unsure{\xi|\zeta}_{\unsure{\sigma|\deleted{\sigma}\sigma}}}\ ,
\]
où
% archive sheet number: 64
